\booksection{Приложение к главе 20: Termux — портативная среда Linux-терминала на устройстве Android}

\chapterintro{Если вам кажется, что режим разработчика недостаточно «хардкорен» и вы хотите опробовать полноценную командную строку Linux на телефоне, то Termux — именно то, что нужно. Без прав root он превращает обычный Android-телефон в портативную Linux-рабочую станцию — мощный инструмент для продвинутых пользователей, исследующих мобильную разработку и администрирование.}

\booksubsection{I. Что такое Termux}

Termux — это приложение с открытым исходным кодом, эмулятор терминала и среда Linux, специально разработанное для устройств Android. Оно предоставляет полноценный опыт Linux shell на телефоне без прав root. Как лёгкое приложение (установочный пакет около 30 МБ), Termux поддерживает богатую экосистему программ через систему управления пакетами APT, позволяя пользователям запускать на мобильном устройстве Python, Node.js, Git, SSH и другие инструменты, и даже развёртывать серверы или проводить тесты на проникновение.

Ключевые преимущества Termux — лёгкость и высокая расширяемость, подходящие для множества сценариев: разработки, обучения и системного администрирования. Он объединяет стандартные пакеты с точной терминальной эмуляцией в одном решении с открытым исходным кодом.

\booksubsection{II. Основные функции и характеристики}

\subsubsection*{Полноценная среда Linux}

Termux\textbf{встроенно включает только Bash}в качестве оболочки. Если вы хотите использовать другие оболочки, такие как Zsh или fish, их нужно установить вручную через менеджер пакетов, например:\texttt{pkg install zsh}или\texttt{pkg install fish}. Termux также поддерживает редакторы Vim, Nano, Emacs, предоставляя почти нативный опыт командной строки. Пользователи могут установить 2000+ пакетов через «pkg install», включая GCC, OpenSSH, SQLMap и другие инструменты разработки и безопасности.

\subsubsection*{Мощное управление пакетами}

Использует менеджер пакетов APT (аналог Debian и Ubuntu), позволяющий установить тысячи Linux-пакетов, включая инструменты разработки, сетевые утилиты и др. Поддерживает среды нескольких языков программирования: Python, Ruby, Node.js, Go, Rust, Swift и другие.

\subsubsection*{Безопасность и сетевые функции}

\begin{itemize}
  \item \textbf{Удалённое SSH-соединение}: поддерживает функции SSH-клиента и сервера, позволяет удалённо подключаться к другим устройствам или превратить телефон в SSH-хост.
  \item \textbf{Инструменты шифрования}: предоставляет OpenSSL и другие инструменты шифрования, поддерживает шифрование данных и безопасную связь.
  \item \textbf{Развёртывание сервера}: можно настроить Nginx или Apache для развёртывания статического веб-сервера.
\end{itemize}

\subsubsection*{Системная интеграция и расширения}

\begin{itemize}
  \item Доступ к разрешениям хранилища телефона для локального управления файлами.
  \item Вызов датчиков, камеры и других аппаратных функций Android через плагин Termux:API.
  \item Поддержка Bluetooth-клавиатур и внешних дисплеев, полноценная поддержка мыши.
\end{itemize}

\booksubsection{III. Экосистема плагинов}

Основное приложение Termux можно расширить с помощью нескольких дополнительных плагинов, предоставляющих дополнительные функции и интеграцию:

\begin{center}
\begin{tabularx}{\linewidth}{|X|X|}
\hline
\textbf{Название плагина} & \textbf{Описание функции} \tabularnewline
\hline
Termux:API & Предоставляет программам командной строки доступ к функциям Android API \tabularnewline
Termux: Boot & Позволяет запускать скрипты при загрузке устройства \tabularnewline
Termux: Float & Предоставляет плавающее окно терминала \tabularnewline
Termux: Styling & Включает цветовые схемы и настройку шрифтов для терминала \tabularnewline
Termux: Tasker & Интегрирует скрипты Termux с приложением автоматизации Tasker \tabularnewline
Termux: Widget & Позволяет запускать скрипты из виджета на главном экране \tabularnewline
\hline
\end{tabularx}
\end{center}

\begin{tcbwarn}{Важное замечание}
Termux и его плагины должны быть установлены из одного и того же источника, поскольку они используют общий sharedUserId. Смешивание установки из разных источников (например, F-Droid и GitHub) приведёт к ошибке установки.

\end{tcbwarn}

\booksubsection{IV. Руководство по установке и настройке}

\subsubsection*{Шаги установки}

\begin{enumerate}
  \item \textbf{Скачайте приложение}: скачайте и установите Termux из официальных источников (F-Droid или GitHub), права root не требуются.
  \item \textbf{Инициализация системы}: после первого запуска дождитесь завершения «Installing bootstrap packages».
  \item \textbf{Получите разрешения хранилища}: выполните\texttt{termux-\allowbreak setup-\allowbreak storage}для получения доступа к файловой системе телефона.
  \item \textbf{Обновите источники ПО}: выполните\texttt{pkg update \&\& pkg upgrade}для обновления списка пакетов.
\end{enumerate}

\subsubsection*{Базовая оптимизация}

\begin{itemize}
  \item \textbf{Смена зеркала}: используйте команду\texttt{termux-\allowbreak change-\allowbreak repo}для смены зеркала через графический интерфейс (рекомендуется зеркало университета Цинхуа).
  \item \textbf{Настройка горячих клавиш}: клавиша уменьшения громкости эмулирует Ctrl, клавиша увеличения громкости — Esc, повышая эффективность ввода.
  \item \textbf{Оформление интерфейса}: установите Zsh + Oh My Zsh, настройте тему и плагины.
\end{itemize}

\booksubsection{V. Типичные сценарии применения}

\subsubsection*{Мобильная разработка и тестирование}

Развёртывание полной среды разработки на телефоне, написание и выполнение скриптов Python, компиляция программ на C с помощью GCC, управление репозиториями кода через Git, синхронизация с удалёнными серверами через SSH.

\subsubsection*{Администрирование серверов}

Удалённое управление серверами через SSH, выполнение системного обслуживания, анализа логов и других задач — решение проблем сервера в любое время и в любом месте.

\subsubsection*{Изучение работы с Linux}

Предоставляет реальную среду Linux — идеальную платформу для изучения командной строки и системного администрирования. Можно практиковать написание Shell-скриптов, управление правами доступа к файлам Linux (например,\texttt{chmod 755}устанавливает биты прав rwxr-xr-x,\texttt{chown user: group}изменяет владельца файла,\texttt{ls -\allowbreak l}показывает три группы по девять битов прав: owner/group/other) и другие базовые команды. Обратите внимание: это классическая Unix/Linux модель прав «три группы по девять битов» (owner / group / other × чтение / запись / выполнение), которая принципиально отличается от списков управления доступом (ACL) Windows NTFS — это две совершенно разные системы.

\subsubsection*{Тестирование сетевой безопасности}

Установка инструментов безопасности Nmap, SQLMap, Hydra, Metasploit для сетевого сканирования и тестирования на проникновение.

\booksubsection{VI. Версия и совместимость}

Последняя стабильная версия Termux — v0.118.2, подходит для устройств с Android $\geq$ 7. Обратите внимание, что Termux может быть нестабилен на Android 12 и выше; рекомендуется следить за официальным журналом обновлений для получения актуальной информации о совместимости.

\begin{tcbinfo}{Резюме}
Termux принёс мощный опыт Linux-терминала на устройства Android, превратив телефон в портативный инструмент разработки и обучения. Разработчики, системные администраторы и технические энтузиасты могут использовать Termux для полноценного опыта командной строки на мобильном устройстве, реализуя концепцию «Linux в кармане».

\end{tcbinfo}

\begin{tcbpractice}{Попробуйте: выполните первую Linux-команду на телефоне}
\begin{enumerate}
  \item Скачайте и установите Termux с F-Droid, при первом запуске дождитесь завершения инициализации.
  \item Выполните\texttt{pkg update \&\& pkg upgrade}для обновления пакетов.
  \item Установите Python:\texttt{pkg install python}.
  \item Введите\texttt{python3 -\allowbreak c "print('Hello from Termux!')"}— ощутите радость выполнения кода на телефоне.
\end{enumerate}

\end{tcbpractice}
